% Part: many-valued-logic
% Chapter: three-valued-logics
% Section: lukasiewicz

\documentclass[../../../include/open-logic-section]{subfiles}

\begin{document}

\olfileid{mvl}{thr}{luk}

\olsection{Logika \L ukasiewicz}

Salah sawijining usulan logika mawa akeh nilai kang paling wiwitan
diterbitake lan diandharake kanthi rinci asalé saka filsuf Polandia,
Jan \L ukasiewicz, ing taun 1921. \L ukasiewicz kabangun dening
masalah perang laut Aristoteles: katon yen \emph{dina iki}, ukara
``Sesuk bakal ana perang laut'' durung bener lan durung salah;
nilai kabenerane durung katetepake. \L ukasiewicz ngusulake
nilai kabeneran katelu kanggo ukara ``kontingen mangsa ngarep''
kaya mangkono.
\begin{quote}
Aku bisa nganggep tanpa kontradiksi yen aku ana ing Warsawa ing
sawijining wektu taun ngarep, umpama awan tanggal 21 Desember,
saiki durung katetepake kanthi positif utawa negatif. Mula bisa
wae, nanging ora mesthi, aku bakal ana ing Warsawa ing wektu
kasebut. Miturut anggepan iki, proposisi ``Aku bakal ana ing
Warsawa nalika awan tanggal 21 Desember taun ngarep'' saiki
durung bisa diarani bener utawa salah. Yen saiki bener, aku ana
ing Warsawa mengko mesthi kudu kelakon, lan iki nalisir
anggepan mau. Dene yen saiki salah, aku ana ing Warsawa mengko
mesthi ora mungkin, lan iki uga nalisir anggepan mau. Mula
proposisi kasebut saiki durung bener lan durung salah, lan
kudu nduweni nilai katelu kang beda saka ``0'' utawa salah,
lan ``1'' utawa bener. Nilai iki bisa kita tandhani nganggo
``$\frac{1}{2}$.'' Nilai iki makili ``kang mungkin'' lan
nggabung karo ``kang bener'' lan ``kang salah'' minangka
nilai katelu.
\end{quote}
Kita nggunakake $\Undef$ kanggo nilai kabeneran katelu
miturut \L ukasiewicz.\footnote{Ing kene \L ukasiewicz
nganggo tembung ``mungkin'' kanthi makna kang saiki ora lumrah,
yaiku mungkin nanging ora mesthi.}

Fungsi kabeneran kanggo konektif $\lnot$, $\land$, lan $\lor$
gampang ditemtokake miturut tafsir iki: negasi ukara
kontingen mangsa ngarep uga kontingen mangsa ngarep, mula
$\tf{\lnot}(\Undef) = \Undef$. Yen salah siji komponen
konjungsi durung katetepake dene liyane bener, konjungsine
uga durung katetepake: gumantung marang kepriye komponen
kontingen iku kelakon, konjungsi bisa dadi bener utawa salah.
Mula \[
    \tf{\land}(\True, \Undef) =
\tf{\land}(\Undef, \True) =
\Undef.
\]
Nanging yen komponen liyane salah, konjungsi ora bisa dadi
bener, mula \[\tf{\land}(\False, \Undef) =
\tf{\land}(\Undef, \False) = \False.\]
Nilai liyane, yen argumene nilai kabeneran kang wis katetepake,
$\True$ utawa $\False$, padha karo ing logika klasik.

Kanggo kondisional, kahanane rada luwih ruwet. Upamakna $q$
ukara kontingen mangsa ngarep. Yen $p$ salah, $p \lif q$
bakal bener ora preduli kepriye $q$ kelakon, mula kita netepake
$\tf{\lif}(\False, \Undef) = \True$. Yen $p$ bener,
$q \lif p$ uga bener ora preduli kepriye $q$ kelakon,
mula $\tf{\lif}(\Undef, \True) = \True$. Yen $p$ bener,
$p \lif q$ bisa dadi bener utawa salah, mula
$\tf{\lif}(\True, \Undef) = \Undef$. Uga, yen $p$ salah,
$q \lif p$ bisa dadi bener utawa salah, mula
$\tf{\lif}(\Undef, \False) = \Undef$. Kari kasus
nalika $p$ lan $q$ loro-lorone kontingen mangsa ngarep.
Miturut motivasi sadurunge, kita sakmesthine menehi nilai
$\Undef$ ing kasus iki. Nanging iki njalari $!A \lif !A$
\emph{dudu} tautologi. \L ukasiewicz ora kabotan ninggalake
$!A \lor \lnot !A$ lan $\lnot(!A \land \lnot !A)$,
nanging ora gelem ninggalake $!A \lif !A$. Mula dheweke
netepake $\tf{\lif}(\Undef, \Undef) = \True$.

\begin{defn}\ollabel{def:lukasiewicz}
Logika telung nilai \L ukasiewicz ditetepake nganggo matriks iki:
\begin{enumerate}
  \item Basa proposisional baku $\Lang L_0$ kang ngemot
  $\lfalse$, $\lnot$, $\land$, $\lor$, $\lif$.
  \item Himpunan nilai kabeneran $V = \{\True, \Undef, \False\}$.
  \item $\True$ mung siji-sijine nilai pinilih, yaiku
  $V^+ = \{\True\}$.
  \item Kanggo tetapan proposisional,
  $\tf{\lfalse}[\LogLuk[3]] = \False$. Fungsi kabeneran liyane
  diwenehake dening tabel iki:
  \begin{center}
    \begin{tabular}{c|c} 
      $\tf{\lnot}$ & \\ 
      \hline  
      $\True$ & $\False$ \\ 
      $\Undef$ & $\Undef$ \\
      $\False$ & $\True$ 
    \end{tabular}
    \quad
    \begin{tabular}{c|ccc} 
      $\tf{\land}[\LogLuk[3]]$ & $\True$ & $\Undef$ & $\False$ \\ 
      \hline 
      $\True$ & $\True$ & $\Undef$ & $\False$ \\ 
      $\Undef$ & $\Undef$ & $\Undef$ & $\False$\\ 
      $\False$ & $\False$ & $\False$ & $\False$ 
    \end{tabular}
    \\[2ex]
    \begin{tabular}{c|ccc} 
      $\tf{\lor}[\LogLuk[3]]$ & $\True$ & $\Undef$ & $\False$ \\ 
      \hline 
      $\True$ & $\True$ & $\True$ & $\True$ \\ 
      $\Undef$ & $\True$ & $\Undef$ & $\Undef$ \\
      $\False$ & $\True$ & $\Undef$ & $\False$ 
    \end{tabular}
    \quad
    \begin{tabular}{c|ccc} 
      $\tf{\lif}[\LogLuk[3]]$ & $\True$ & $\Undef$ & $\False$ \\ 
      \hline 
      $\True$ & $\True$ & $\Undef$ & $\False$ \\ 
      $\Undef$ & $\True$ & $\True$ & $\Undef$  \\ 
      $\False$ & $\True$ & $\True$ & $\True$ 
    \end{tabular}
  \end{center} 
\end{enumerate}
\end{defn}

Kaya kang gampang dideleng, saben !!{formula}~$!A$ kang mung
ngemot $\lnot$, $\land$, lan $\lor$ bakal nduweni nilai
kabeneran~$\Undef$ yen kabeh !!{propositional variable}s
diparingi nilai~$\Undef$. Mula, umpama, tautologi klasik
$p \lor \lnot p$ lan $\lnot(p \land \lnot p)$ dudu
tautologi ing $\LogLuk[3]$, amarga $\pValue{v}(!A) = \Undef$
nalika $\pAssign v(p) = \Undef$.

Ing !!{valuation}s kang nyukupi $\pAssign v(p) = \True$ utawa
$\False$, $\pValue v(!A)$ padha karo nilai kabeneran klasik.

\begin{prop}
  Yen $\pAssign v(p) \in \{\True, \False\}$ kanggo saben $p$
  ing~$!A$, mula
  $\pValue v(!A)[\LogLuk[3]] = \pValue v(!A)[\LogCL]$.
\end{prop}

\begin{prob}\label{mvl:thr:luk:prob:luk-iff} Upamakna kita netepake
  $\pValue v(!A \liff !B) = \pValue v((!A \lif !B) \land (!B \lif
  !A))$ ing~$\LogLuk[3]$. Kaya apa tabel kabeneran kanggo $\liff$?
\end{prob}

Akeh tautologi klasik \emph{uga} tautologi ing \LogLuk[3],
umpamane $\lnot p \lif (p \lif q)$. Kaya ing logika klasik,
kita bisa mriksa kanthi tabel kabeneran:
\begin{center}
\begin{tabular}{cc|cccccc}
  $p$ & $q$ & $\lnot$ & $p$ & $\lif$ & $\smash{(}p$ & $\lif$ & $q\smash{)}$ \\ \hline
  \True & \True & \False & \True & \True & \True & \True & \True \\
  \True & \Undef & \False & \True & \True & \True & \Undef & \Undef \\
  \True & \False & \False & \True & \True & \True & \False & \False \\
  \Undef & \True & \Undef & \Undef & \True & \Undef & \True & \True \\
  \Undef & \Undef & \Undef & \Undef & \True & \Undef & \True & \Undef \\
  \Undef & \False & \Undef & \Undef & \True & \Undef & \Undef & \False \\
  \False & \True & \True & \False & \True & \False & \True & \True \\
  \False & \Undef & \True & \False & \True & \False & \True & \Undef \\
  \False & \False & \True & \False & \True & \False & \True & \False \\  
\end{tabular}
\end{center}

\begin{prob}
  Tuduhna yen formula-formula iki tautologi ing \LogLuk[3]:
  \begin{enumerate}
    \item $p \lif (q \lif p)$
    \item\label{mvl:thr:luk:prob:luk-taut-2} $\lnot(p \land q) \liff (\lnot p \lor \lnot q)$
    \item\label{mvl:thr:luk:prob:luk-taut-3} $\lnot(p \lor q) \liff (\lnot p \land \lnot q)$
  \end{enumerate}
  (Ing \olref[mvl][thr][luk]{prob:luk-taut-2} lan
  \olref[mvl][thr][luk]{prob:luk-taut-3}, gunakna $!A \liff !B$
  minangka cekakan kanggo $(!A \lif !B) \land (!B \lif !A)$,
  utawa delengen jawabanmu marang
  \olref[mvl][thr][luk]{prob:luk-iff}.)
\end{prob}

\begin{prob}
  Tuduhna yen tautologi-tautologi klasik iki dudu tautologi ing~\LogLuk[3]:
  \begin{enumerate}
    \item $((\lnot p \land p) \lif q)$
    \item $((p \lif q) \lif p) \lif p$
    \item $(p \lif (p \lif q)) \lif (p \lif q)$
  \end{enumerate}
\end{prob}

Saka kene bisa wae dikira yen sanajan ora kabeh tautologi
klasik uga tautologi ing~$\LogLuk[3]$, paling ora
nilainé~$\True$ utawa~$\Undef$ kanggo saben
!!{valuation}. Nanging ora mangkono. Kontraconto diwenehake dening
\[
  \lnot(p \lif \lnot p) \lor \lnot(\lnot p \lif p)
\]
kang nilaine $\False$ yen $p$ iku~$\Undef$.

\begin{prob}
  Endi sesambungan ing ngisor iki kang lumaku ing logika \L ukasiewicz?
  Wenehana tabel kabeneran kanggo saben sesambungan.
  \begin{enumerate}
    \item $p, p \lif q \Entails q$
    \item $\lnot\lnot p \Entails p$
    \item $p \land q \Entails p$
    \item $p \Entails p \land p$
    \item $p \Entails p \lor q$
  \end{enumerate}
\end{prob}

\L ukasiewicz ngarep-arep bisa mbangun logika kemungkinan adhedhasar
sistem telung nilaine, kanthi nambah konektif sapanggonan $\Diamond
!A$ (tegesé ``$!A$ mungkin'') lan $\Box !A$ kang cocog (tegesé ``$!A$
mesthi''):
\begin{center}
  \begin{tabular}{c|c} 
    $\tf{\Diamond}$ & \\ 
    \hline  
    $\True$ & $\True$ \\ 
    $\Undef$ & $\True$ \\
    $\False$ & $\False$ 
  \end{tabular}
  \quad
  \begin{tabular}{c|c} 
    $\tf{\Box}$ & \\ 
    \hline  
    $\True$ & $\True$ \\ 
    $\Undef$ & $\False$ \\
    $\False$ & $\False$ 
  \end{tabular}
\end{center}
Tegese, $p$ mungkin yen lan mung yen durung katetepake salah;
dene $p$ mesthi yen lan mung yen wis katetepake bener.

\begin{prob}
  Tuduhna yen $\Box p \liff \lnot\Diamond \lnot p$ lan $\Diamond p \liff
  \lnot \Box \lnot p$ iku tautologi ing~\LogLuk[3] kang ditambahi
  tabel kabeneran kanggo $\Box$ lan~$\Diamond$.
\end{prob}

Nanging, kekurangan usulan logika modal iki enggal katon:
kepiye wae kahanane mengko, $p \land \lnot p$ ora bakal bener.
Mula sanajan saiki durung dipesthekake (lan mulane nilai kabenerane durung katetepake),
ukara iku kudu dianggep ora mungkin; tegesé,
$\lnot \Diamond(p \land \lnot p)$ sakmesthine tautologi.
Nanging yen $\pAssign v(p) = \Undef$, mula
$\pValue v(\lnot \Diamond(p \land \lnot p)) = \False$.
Senajan \L ukasiewicz bener yen rong nilai kabeneran ora cukup
kanggo mbedakake kemungkinan lan kepastian, nambah nilai katelu
uga isih durung cukup.

\end{document}
